%------------------------------------------------------------------------------%
\section{Generalização da Noção de Ângulo}

Após uma volta completa, continuamos medindo os ângulos. Assim, por exemplo, o
ângulo $\sfrac{\pi}{6}$ ocupa a mesma posição que os ângulos
$2\pi + \sfrac{\pi}{6}$
(após uma volta) ou
$4\pi+\sfrac{\pi}{6}$ (após duas voltas) e assim por diante. Mais geralmente,
para qualquer $n$ natural, os ângulos $x$ e $x+2\pi n$ ocupam a mesma posição
no círculo trigonométrico.

Para dar sentido a ângulos negativos, convenciona-se que ao girarmos no sentido
anti-horário, medimos os ângulos negativamente.

%------------------------------------------------------------------------------%
\begin{figure}
  \centering
  \includegraphics[width=0.3\textwidth,trim={0 0 10,5cm 1cm},clip]{circ3}
\end{figure}
%------------------------------------------------------------------------------%

Portanto, um ângulo $x$ radianos ocupa o mesmo lugar que todos os ângulos da forma
\[
  x+2\pi n, \quad n\in\Z.
\]

%------------------------------------------------------------------------------%